\documentclass[10pt,journal,compsoc]{IEEEtran}
\usepackage[utf8]{inputenc}
\usepackage{amsmath,amssymb,amsfonts}
\usepackage{graphicx}
\usepackage{cite}
\usepackage{booktabs}
\usepackage{hyperref}

\begin{document}

\title{Hierarchical Transformer Architectures for Multilead Biomedical Signal Localization}

\author{Daniandra Prayudisty,~\IEEEmembership{Member,~IEEE}
\thanks{Daniandra Prayudisty is with Yudiaz Creative Studio and School of Computing, Telkom University, Bandung, Indonesia.}
}

\markboth{Journal of Biomedical and Health Informatics,~Vol.~XX, No.~X, \today}%
{Prayudisty: Hierarchical Transformer Architectures}

\IEEEtitleabstractindextext{%
\begin{abstract}
Precise anatomical localization of focal cardiac arrhythmias from standard 12-lead electrocardiograms (ECG) plays a critical role in pre-procedural planning for radiofrequency catheter ablation. In this work, we propose a leak-free, patient-wise validated deep learning framework integrating convolutional feature representation with hierarchical transformer networks. Experimental evaluation on patient-held-out validation demonstrates robust macro-F1 improvements across anatomical subregions.
\end{abstract}

\begin{IEEEkeywords}
Arrhythmia localization, 12-lead ECG, Transformer, Catheter ablation, Patient-wise validation.
\end{IEEEkeywords}}

\maketitle
\IEEEdisplaynontitleabstractindextext
\IEEEpeerreviewmaketitle

\section{Introduction}
\IEEEPARstart{F}{ocal} ventricular arrhythmias originating from outflow tracts present diagnostic challenges in non-invasive clinical mapping. Traditional manual algorithmic criteria suffer from inter-observer variability and limited generalization across heterogeneous cardiac morphologies.

\section{Methodology}
\subsection{Signal Preprocessing}
Standard 12-lead ECG records are sampled and denoised using discrete wavelet transform (DWT). Signal normalization is conducted strictly within inner training folds to eliminate data leakage.

\begin{equation}
y(t) = \sum_{k} c_k \psi_{j,k}(t) + \sum_{k} d_k \phi_{j,k}(t)
\end{equation}

\subsection{Architecture Overview}
The network comprises spatial-temporal convolutional feature extractors followed by hierarchical self-attention heads:
\begin{equation}
\text{Attention}(Q, K, V) = \text{softmax}\left(\frac{QK^T}{\sqrt{d_k}}\right)V
\end{equation}

\section{Results and Discussion}
Table~\ref{tab:results} summarizes validation metrics across patient-wise stratified evaluations.

\begin{table}[h]
\caption{Cross-Validation Performance Comparison}
\label{tab:results}
\centering
\begin{tabular}{lccc}
\toprule
\textbf{Model Architecture} & \textbf{Accuracy} & \textbf{Macro-F1} & \textbf{AUC} \\ \midrule
Baseline CNN & 64.2\% & 63.8\% & 0.741 \\
Transformer Encoder & 68.5\% & 67.9\% & 0.785 \\
Hierarchical Transformer & \textbf{72.5\%} & \textbf{72.1\%} & \textbf{0.824} \\ \bottomrule
\end{tabular}
\end{table}

\section{Conclusion}
The proposed hierarchical approach provides high fidelity localization while preserving strict patient-level boundary validation.

\begin{thebibliography}{00}
\bibitem{ref1} D. Prayudisty, ``Deep Learning Frameworks for Arrhythmia Localization,'' \emph{IEEE Trans. Biomed. Eng.}, vol. 71, pp. 110--122, 2026.
\end{thebibliography}

\end{document}
